\chapter{A27 电话} % Introduction chapter suppressed from the table of contents

\hypertarget{ux66feux7ecfux4fc3ux8fdbux4ea4ux6d41ux7684ux7535ux8bddux73b0ux5728ux5374ux65e5ux76caux6210ux4e3aux4ea4ux6d41ux7684ux969cux788d}{%
\section{曾经促进交流的电话，现在却日益成为交流的障碍}\label{ux66feux7ecfux4fc3ux8fdbux4ea4ux6d41ux7684ux7535ux8bddux73b0ux5728ux5374ux65e5ux76caux6210ux4e3aux4ea4ux6d41ux7684ux969cux788d}}

(56 The telephone, which once facilitated communication, now
increasingly obstructs it.)

\framebox{%
\begin{minipage}[t]{0.97\columnwidth}\raggedright
Prof Ackoff 经验之谈：\\
如今，接电话的人可以这样巧妙地告诉打电话的人不想被联系：先将电话置于等待状态；最后再挂掉。

过去，只需要拨打七位数字即可联系到大多数的人（美国本地电话仅需七位数字）。

如今，电话号码的位数和等待时间成倍增加，远超常人的记忆力和耐心所能承受的范围。

原本希望通过电话找到合适的联系人，却常常因人工生成的电话转接号码无效而徒劳无功。

即便偶尔联系到某人，却往往因为对方没有能力响应需求，也不知应转接何人，又引发一连串的转接，最终以精心设计的``意外''中断告终。

电话曾帮我们大幅减少书面沟通的需求，如今却迫使我们重拾传统的沟通方式。
\strut
\end{minipage}}

\textbf{严重故障需求流程图} （图27）\\
%\href{文件:value_1-3.jpg}{600px}

\includegraphics[width=6cm]{value_1-3.jpg}

\newpage

不知道大家每天会收到多少骚扰电话？通常，我的反应是直接问``哪位？''，确认是陌生人便直接挂断。\\
在教授所处的年代，骚扰电话尚未泛滥，被呼叫者还能礼貌拒绝。\\
如今，我倒有点同情那些被迫拨打骚扰电话的员工，他们每日重复这些工作，想必备受煎熬。问题的根源在于管理层，他们误以为电话促销仍有效，殊不知这种方式不仅效果近乎为零，还损害公司形象。在互联网时代，搜索引擎的功能强大如此，如果产品足够优秀，只需在网站展示并得到用户的点赞，客户自会主动搜索，无需电话推销。\\
正因电话的滥用，人们更渴望免受打扰，可以安静地专注于自己的工作。\\

每次听到``查询余额按1，转账按2，优惠活动按3\ldots{}\ldots{}''这样的自动应答系统录音，我就会直接挂断。经验告诉我，指望这些系统永远无法连通我需要的服务。不少银行、售票等行业普遍采用此类系统一些有经验的客户经理会直接告知正确按键顺序（如1342）或人工客服接入方式，以节省时间。或如何可以直接与人工客服对话，减小对我的耽误。如今，随着科技不断进步，部分公司引入机器人模拟真人对话，但效果未必好。如某售票系统的机器人会热情地问候：``早上好！小倩很高兴为您服务，请问有什么可以帮助您的？''\\
然而，两分钟的``对话''后，系统又提示输入身份证的后四位号码，而我只有港澳通行证号，最终还是无法购票。（但若是与真人通话，整个交易通常在2-3分钟即可完成。）企业误以为机器人（高科技）能提升客户满意度，实则适得其反，客户只会更觉烦躁。所以我总想绕过系统自动应答，直达人工客服，却从未成功破解其设计。\\
在精益价值流培训过程中，我会展示前面的图27，并提问：``如果现在处理存在严重缺陷的流程，从提出申请到解决需要10天。请问有何建议可以缩短时间、减少浪费的？''\\
学员们经过讨论并提交建议后，我接着再问：``是否考虑让用户直接联系开发工程师？''\\
他们又回应：``开发和测试都太忙了，如无客服预先过滤，工程师无法承受这么大的工作量。''\\
接着我分享如下案例：``某公司有800位工程师，在传统客服过滤的模式下，工程部用于处理严重缺陷约占60\%的工作量。变更流程后，通过工程师直接接听客户电话，工作量降至24\%。''
学生觉得不可思议，我接着说：``通过工程师接听客户电话，他们同步提供了客户缺陷与对应的解决办法，并更新到知识库，供客户自查。这不仅将响应时间缩短至原来的三分之一，还让工程师直接收集到用户使用系统过程中遇到的问题，深入了解客户的使用场景，为下一版本的改进提供洞见。''\\
唯一不变的唯有改变。 (The only thing that is constant is change / Change
is the only constant in life)\\
若能跳出思维定式，便能海阔天空，激发无限创意。


